\documentclass[11pt]{article}
\usepackage[T1]{fontenc}
\usepackage{lmodern,amsmath,amssymb,amsthm,microtype}
\usepackage[margin=1in]{geometry}
\usepackage[colorlinks=true,linkcolor=blue,urlcolor=blue,citecolor=blue]{hyperref}
\usepackage{enumitem}
\setlist{nosep}
\newtheorem{theorem}{Theorem}
\newtheorem{corollary}[theorem]{Corollary}
\newtheorem{proposition}[theorem]{Proposition}
\newtheorem{lemma}[theorem]{Lemma}
\newcommand{\rev}{\operatorname{rev}}
\newcommand{\len}{\ell}
\newcommand{\ord}{\operatorname{ord}}
\title{Exact-length palindrome residues and limits of\newline arithmetic invariants for reverse-and-add}
\author{Research note}
\date{13 October 2026}
\begin{document}
\maketitle
\begin{abstract}
For every fixed positive modulus, we construct decimal palindromes of
every sufficiently large admissible length in exactly the residue classes
allowed by the units digit and the even-length divisibility rule. An explicit
elementary bound handles arbitrary factors of the modulus. It follows that
no nonempty palindrome-free forward invariant of decimal reverse-and-add
can have membership determined solely by exact digit length and a fixed
congruence class. We also classify the moduli for which reversal has an
autonomous residue update and give a direct obstruction to a continuous
10-adic extension. These are limitations of particular proof methods.
\textbf{They do not prove or disprove the existence of a decimal Lychrel
number.} No publication-priority claim is made.
\end{abstract}

\section{The exact-length residue theorem}
For a positive integer $n$, let $\len(n)$ be its ordinary decimal digit
length, let $\rev(n)$ be its decimal reversal with leading zeros discarded,
and set $T(n)=n+\rev(n)$. A palindrome satisfies $\rev(n)=n$.
All numbers and all words in the constructions below are finite.

Fix $M\geq1$ and write uniquely
\[
 M=uq,\qquad u=2^\alpha5^\beta,\qquad \gcd(q,10)=1.
\]
Put $k=\max(1,\alpha,\beta)$, and let $r=\ord_q(10)$ when $q>1$;
use the convention $r=1$ when $q=1$. Define
\[
 B(M)=2k+4qr.
\]

\begin{theorem}\label{thm:saturation}
For every integer $L\geq B(M)$ and every residue $a$ modulo $M$, there is
a positive $L$-digit decimal palindrome $P\equiv a\pmod M$ if and only if
\begin{enumerate}[label=(\roman*)]
\item $10$ does not divide both $M$ and $a$;
\item when $L$ is even, $\gcd(M,11)$ divides $a$.
\end{enumerate}
The same bound works for every larger length, without a restriction on
$L$ modulo $r$.
\end{theorem}

\begin{proof}
Necessity follows because a positive palindrome cannot end in zero, and
an even-length decimal palindrome is divisible by $11$.

First choose $A$ such that
\[
 0\leq A<10^k,\qquad A\equiv a\pmod u,\qquad 10\nmid A.
\]
Start with the least nonnegative residue $a\bmod u$. If it ends in zero,
add $u$. If $10\mid u$, condition (i) makes that addition unnecessary.
Otherwise $u$ is $1$, a power of $2$, or a power of $5$. In that case
$2u\leq10^k$, and adding $u$ to a residue below $u$ both stays below
$10^k$ and makes its units digit nonzero.

Write $A$ as a zero-padded block of $k$ digits. Place it at the low end of
an $L$-digit word and place its mirror at the high end, filling the middle
with zeros. The resulting palindrome $C$ has exactly $L$ digits because
$A$ has nonzero units digit. Since $u\mid10^k$,
$C\equiv a\pmod u$. If $L$ is even, $11\mid C$.

Number digit positions from $0$ at the units end. A digit $1$ placed at
both positions $i$ and $L-1-i$ contributes
\[
 c_i=10^i+10^{L-1-i}.
\]
Consider the two disjoint intervals of lower-half positions
\[
 I_0=[k,k+qr-1],\qquad I_1=[k+qr,k+2qr-1].
\]
Their largest position is at most $\lfloor(L-2)/2\rfloor$. Consequently
the corresponding mirror pairs are mutually disjoint and avoid the
already fixed outer blocks. Exactly $q$ positions in $I_0$ are congruent
to $0$ modulo $r$, and exactly $q$ positions in $I_1$ are congruent to $1$
modulo $r$. For $r=1$ the residue labels coincide, but the two pools remain
disjoint.

Modulo $q$, the weights in these respective pools are
\[
 c_0=1+10^{L-1},\qquad c_1=10+10^{L-2}.
\]
These unreduced integers satisfy
\[
 10c_1-c_0=99.
\]
Thus $g=\gcd(q,c_0,c_1)$ divides $99$. Since $c_0\equiv2\pmod3$,
$g$ has no factor $3$, so $g\mid11$. Modulo $11$, both weights vanish
for even $L$, whereas $c_0\equiv2$ for odd $L$. Therefore
\[
 g=\begin{cases}\gcd(q,11),&L\text{ even},\\1,&L\text{ odd}.\end{cases}
\]
This argument also handles arbitrary powers of $11$ in $q$: only one
factor of $11$ can divide the common gcd.

If $q=1$, the already constructed $C$ completes the proof. Otherwise,
the subgroup generated by $c_0,c_1$ in $\mathbb Z/q\mathbb Z$ consists
exactly of the multiples of $g$. The hypotheses and the choice of $C$
give $g\mid(a-C)$. Choose $x,y\in\{0,\ldots,q-1\}$ satisfying
\[
 xc_0+yc_1\equiv a-C\pmod q.
\]
Place a digit $1$ in any $x$ mirror pairs from the first pool and any $y$
from the second, leaving other free digits zero. This creates an $L$-digit
palindrome $P\equiv a\pmod q$. Every new position is at least $k$,
so $P\equiv C\equiv a\pmod u$ still holds. Since $\gcd(u,q)=1$,
$P\equiv a\pmod M$.
\end{proof}

\section{A consequence for invariant-set proofs}
An invariant proof of a Lychrel witness would exhibit a nonempty set
$S$ of positive integers with $T(S)\subseteq S$ and with no palindrome
in $S$. The next result excludes one broad, precisely specified way of
describing such a set. It does not exclude invariants using digit patterns,
carries, or other information.

\begin{lemma}\label{lem:orbit}
Every positive reverse-and-add orbit has unbounded digit length and is
divisible by $11$ at time $6$ and every later time. The uniform bound $6$
is sharp. Two consecutive terms cannot both end in zero.
\end{lemma}
\begin{proof}
Since $\rev(n)>0$, the positive integer orbit strictly increases and is
unbounded. An $L$-digit input and its reversal are both below $10^L$, so
the next length is at most $L+1$; the length cannot decrease.

There must be a length increase within five steps. To see this, let $a$
be the sum of the leading and units digits of a current input. It is
at least $1$. In a step without overflow, the new units digit is $a$,
and the new leading digit is $a+c$, where $c\in\{0,1\}$ is the incoming
carry to the leading column. Thus the new outer digit sum is at least
$2a$. Four consecutive steps without overflow would produce an outer
sum at least $16$, which forces overflow at the fifth step. If the
initial length is odd, the first increase therefore reaches an even
length by time $5$; if it is even, no wait is needed. Reversing an
$L$-digit expansion gives
\[
 \rev(n)\equiv(-1)^{L-1}n\pmod {11}.
\]
An even-length input therefore produces a multiple of $11$, and this
divisibility persists at every subsequent step. Finally, if $n$ ends in
zero, the units digit of $T(n)$ is the nonzero leading digit of $n$.
Sharpness follows from the continued reverse-and-add orbit
\[
100,\ 101,\ 202,\ 404,\ 808,\ 1616,\ 7777,
\]
whose first multiple of $11$ is its term at time $6$. Iteration here
continues even when an earlier term is a palindrome, as required by the
definition of $T^t$.
\end{proof}

\begin{corollary}\label{cor:invariant}
Fix $M\geq1$. There is no nonempty set $S\subseteq\mathbb N_{>0}$
having all three properties:
\begin{enumerate}[label=(\roman*)]
\item for positive $n,m$, equality of $\len(n),\len(m)$ and of $n,m$
      modulo $M$ implies $n\in S\Longleftrightarrow m\in S$;
\item $T(S)\subseteq S$;
\item $S$ contains no palindrome.
\end{enumerate}
The allowed residues may depend arbitrarily on the exact length.
\end{corollary}
\begin{proof}
Start an orbit in $S$. By Lemma~\ref{lem:orbit}, choose a late orbit
term $n$ with $\len(n)\geq B(M)$, $11\mid n$, and $10\nmid n$.
Theorem~\ref{thm:saturation} supplies a palindrome of exactly the same
length and the same residue modulo $M$. Property (i) puts that palindrome
in $S$, contradicting (iii).
\end{proof}

This includes a fixed finite list of moduli, by taking their least common
multiple, and a fixed finite list of truncated prime valuations, by taking
the appropriate product of prime powers. It does not imply that the
actual orbit ever visits the constructed palindrome.

There is also a necessary condition for descriptions using a
length-dependent modulus $M(L)$. In any nonempty safe invariant of that
form, every sufficiently late orbit term $n$ with nonzero units digit must
satisfy
\[
 \len(n)<B\bigl(M(\len(n))\bigr).
\]
Otherwise the same contradiction applies. This condition alone supplies
neither an invariant nor a finite Lychrel witness.

\section{Which moduli give an autonomous residue update?}
Even before seeking an invariant, one might try to compute reversal
modulo $M$ from the current residue and length. The possible moduli
admit a complete elementary classification.

\begin{theorem}\label{thm:99}
For a fixed $M\geq1$, the residue $\rev(n)\bmod M$ is determined by
$(\len(n),n\bmod M)$ for all positive integers $n$ if and only if
$M\mid99$. The necessity still holds if this requirement is imposed
only at every sufficiently large length. The same classification holds
for $T(n)\bmod M$.
\end{theorem}
\begin{proof}
If $M\mid99$, then $10^2\equiv1\pmod M$. Termwise reversal gives
\[
 \rev(n)\equiv10^{L-1}n\pmod M,\qquad L=\len(n),
\]
so $T(n)\equiv2n$ for odd $L$ and $T(n)\equiv11n$ for even $L$.

Conversely, if $p\in\{2,5\}$ divides $M$, choose $L$ sufficiently large
that each interval of $L$-digit numbers with leading digit $1$ or $2$
contains a representative of every residue class modulo $M$. Choose two
such representatives of the same class, one in each interval. Their
reversals have distinct residues $1$ and $2$ modulo $p$, a contradiction.
Thus $\gcd(M,10)=1$ is necessary.

Choose $L\geq3$ so large that
$n=10^{L-1}+9$, $n+1$, $n+M$, and $n+M+1$ all have length $L$.
The increment at $n$ carries through exactly one trailing nine. Since
$10\nmid M$, the number $n+M$ does not end in nine, so its increment has
no carry. Hence
\begin{align*}
 \rev(n+1)-\rev(n)&=10^{L-2}-9\cdot10^{L-1},\\
 \rev(n+M+1)-\rev(n+M)&=10^{L-1}.
\end{align*}
If equal input residues at equal lengths determine equal reversal
residues, these two differences agree modulo $M$. Therefore
$M\mid99\cdot10^{L-2}$, and coprimality forces $M\mid99$.
Finally, $T(n)-n=\rev(n)$ makes the two classifications equivalent.
\end{proof}

\section{A discontinuity at every finite positive input}
The completion in powers of $10$ cannot by itself provide a continuous
extension of ordinary reverse-and-add.

\begin{proposition}\label{prop:adic}
No continuous self-map of the $10$-adic integers agrees with $T$ on all
positive ordinary integers.
\end{proposition}
\begin{proof}
Fix $n>0$ of length $L$ and a digit $c\in\{1,\ldots,9\}$. For $j\geq L$,
put $n_j=n+c10^j$. Direct reversal of its digit string gives
\[
 \rev(n_j)=c+10^{j-L+1}\rev(n).
\]
Thus $n_j\longrightarrow n$ $10$-adically, but
$T(n_j)\longrightarrow n+c$. Choose $c\not\equiv\rev(n)\pmod {10}$;
then $n+c\neq T(n)$, contradicting continuity at $n$.
\end{proof}
For example, for every $j\geq3$,
\[
 T(196+4\cdot10^j)=200+1091\cdot10^{j-2}.
\]
The inputs converge to $196$ and the outputs to $200$, whereas
$T(196)=887$. A different state space retaining both ends of a digit
string is not excluded. Any such argument still needs an explicit
finite-seed realization.

\section{Prior work, verification, and remaining gap}
Harminc and Sot\'{a}k~\cite{HS} prove that a nonconstant positive arithmetic
progression contains infinitely many decimal palindromes unless all of
its terms are divisible by $10$. This already rules out a safe invariant
depending on a fixed residue alone. Banks, Hart, and Sakata~\cite{BHS},
Proposition~4.2, give exact-length equidistribution when
$\gcd(M,990)=1$; eventual exact-length support in that case is established
prior work. Col~\cite{Col} strengthens distribution results with the same
coprimality condition. The elementary construction here handles the
support question for arbitrary $M$, with the stated explicit bound.
The limited literature review did not establish whether that full
formulation has appeared before. No novelty or publication-readiness
claim follows from the present review.

The construction was implemented using literal finite digit strings.
Recorded checks constructed $2,512$ compatible cases, including moduli
with powers of $11$ and lengths up to $1,288,411$ digits. An independent
dynamic program computed exact palindrome residue sets for all
$1\leq M\leq300$ and $1\leq L\leq80$. Separate finite controls checked
the residue-update counterexamples. These checks support the implementation;
the arbitrary-length conclusions follow from the proofs above.

\paragraph{Formalization boundary.}
The all-moduli exact-length construction and its invariant-set corollary
have mathematical proofs and an independent proof audit; they are not
claimed as fully formalized Lean theorems. Lean~4.19 with only the standard
library verifies the sharp time-$6$ divisibility theorem, exact reversal
and step recurrences modulo $99$ and every divisor of $99$, the absence
of consecutive zero units digits, and the existence of arbitrarily long
actual orbit states meeting both residue compatibility conditions.
The modular module contains $25$ completed theorems. It also proves that
palindromic overflow from even input width requires $121\mid T(n)$.
These proofs are integrated into the standalone
\texttt{Lychrel196Research.lean} under \texttt{AstraModular}.
Neither the necessity half of Theorem~\ref{thm:99} nor
Proposition~\ref{prop:adic} is claimed as fully formalized here.
The archive distinguishes kernel-checked lemmas, paper proofs, and finite
computations; no unproved research assertion is inserted as a Lean axiom.


For decimal existence the missing statement remains
\[
 \exists n>0\ \ \forall t\geq0:\quad
 \rev\bigl(T^t(n)\bigr)\neq T^t(n).
\]
None of the theorems in this note establishes that statement or its
negation. In particular, constructing seeds for arbitrarily large finite
avoidance horizons would not supply one seed for all horizons.

\begin{thebibliography}{9}
\bibitem{HS} M. Harminc and R. Sot\'{a}k,
\emph{Palindromic Numbers in Arithmetic Progressions},
Fibonacci Quarterly 36(3) (1998), 259--262.
\url{https://www.fq.math.ca/Scanned/36-3/harminc.pdf}.
\bibitem{BHS} W. D. Banks, D. N. Hart, and M. Sakata,
\emph{Almost All Palindromes Are Composite},
Mathematical Research Letters 11 (2004), 853--868.
\url{https://arxiv.org/abs/math/0405056}.
\bibitem{Col} S. Col,
\emph{Palindromes dans les progressions arithm\'{e}tiques},
Acta Arithmetica 137 (2009), 1--41.
\url{https://doi.org/10.4064/aa137-1-1}.
\end{thebibliography}
\end{document}
